\label{sec:res_fem_verification}
Verification of the simulation results was carried out by comparison with the established
analytical solutions for the drawing stress, those of Siebel and of Avitzur, and by internal
consistency tests of the exported fields.
% NOTE:, K18: the section no longer opens with "in place of experimental validation"; it opens
% with the analytical comparison the review asks for.

The drawing stress cannot be read from the cases used for the surrogates, whose exported frame is
taken after the wire has left the die and the section is unloaded (median
$|\int\sigma_{zz}\,\mathrm{d}A|/A$ of 5.3~MPa against an rms $\sigma_{zz}$ of 167~MPa).
% src: VALIDATION.md §1 (5.3 / 167 MPa on the 1756 retained cases)
It is obtained instead from a dedicated export, in which the reaction at the die reference point,
the temperature and the equivalent plastic strain are written as histories in time for four runs
at a reduction of 1.6\,\% in diameter (3.2\,\% in area) and a land ratio of unity. The axial reaction is averaged over the
steady phase of the pass, defined by a die travel of 20 to 100~mm; the signal there is not a
plateau but a stick--slip sawtooth whose relative spread is 13.6\,\% at $\mu = 0.025$ and
4.5\,\% at $\mu = 0.1$, so the time average is used, the quantity an energy method estimates.
Divided by the area actually reached in the model, 985~mm$^2$ at an outlet radius of
17.71~mm, this gives a drawing stress of 54.1, 62.0 and 78.3~MPa at $\alpha = 12^\circ$ for
$\mu = 0.025$, 0.05 and 0.1, and 62.3~MPa at $\alpha = 8^\circ$.
% src: publication/reexport/drawing_force.py on publication/reexport/history/; the median over the
% same window would give 59.2 / 65.5 / 77.6 / 62.4 MPa, above the time mean at mu = 0.025 and 0.05
% and below it at mu = 0.1, which is the sawtooth changing shape with friction.

Two of the three relations compared below are classical. In the slab-method form of Siebel,
\begin{equation}
\sigma_d = \sigma_f\left[\ln\frac{1}{1-r}\,(1+\mu\cot\alpha) + \frac{2}{3}\,\alpha\right],
\label{eq:siebel}
\end{equation}
with $r = 1-(R_f/R_0)^2$ the reduction in area, $\alpha$ the die semi-angle and $\sigma_f$ the flow
stress, and, in the upper-bound form of Avitzur,
\begin{equation}
\sigma_d = \sigma_f\left[\,2f(\alpha)\ln\frac{R_0}{R_f}
  + \frac{2}{\sqrt3}\,\frac{\alpha-\sin\alpha\cos\alpha}{\sin^2\alpha}
  + \frac{2}{\sqrt3}\,m\cot\alpha\ln\frac{R_0}{R_f}
  + \frac{2}{\sqrt3}\,m\,\frac{L}{R_f}\right],
\label{eq:avitzur}
\end{equation}
whose four terms are the ideal deformation work, the redundant shear, the friction on the cone and
the friction on the land; $f(\alpha)$ is the closed form of the deformation integral (verified
against direct numerical integration to eight digits), $f\to1$ as $\alpha\to0$, so that the first
term reduces to $\ln[1/(1-r)]$, $m$ is the friction factor and $L$ the land length
\citep{siebel1947drahtziehen,avitzur1968metal,avitzur1971study}.

Equation~\eqref{eq:avitzur} measures the friction through the factor $m$, the ratio of the
interface shear stress to the shear yield stress, whereas the simulations are driven by a Coulomb
coefficient $\mu$; the Siebel relation above already takes $\mu$ directly. The conversion
$m = \sqrt3\,\mu$ holds only where the normal pressure on the contact equals the flow stress, so it
is an assumption about the contact rather than a change of units. Avitzur gives the bound in a
second form in which $\mu$ enters directly and the land is treated as a bearing along which the
axial stress grows under Coulomb friction \citep{avitzur1968metal,avitzur1971study}; in the writing
of Vega and co-workers \citep{vega2009temperature,vega2009investigation,haddi2011analysis}, who
build their modification on it, it reads
\begin{equation}
\frac{\sigma_d}{\sigma_f} = \frac{\displaystyle
  2f(\alpha)\ln\frac{R_0}{R_f}
  + \frac{2}{\sqrt3}\left(\frac{\alpha}{\sin^2\alpha}-\cot\alpha\right)
  + 2\mu\left[\cot\alpha\left(1-\ln\frac{R_0}{R_f}\right)\ln\frac{R_0}{R_f}
  + \frac{L}{R_f}\right]}
  {\displaystyle 1+2\mu\,\frac{L}{R_f}},
\label{eq:avitzur-vega}
\end{equation}
written here for zero back tension and with the flow stress in place of the yield stress of the
original. The denominator is not a correction: the axial balance on a cylindrical land,
$\mathrm{d}\sigma_z/\mathrm{d}z = 2\mu p/R_f$ closed by the yield condition $p = \sigma_f-\sigma_z$,
integrates to $\sigma_z(L) = \sigma_f-[\sigma_f-\sigma_z(0)]\exp(-2\mu L/R_f)$, of which
Eq.~\eqref{eq:avitzur-vega} is the linearisation; the exact form lies above it by at most
4.0\,\% at the present parameters, so the linearised relation, if anything, understates the
bearing term it assumes. Equation~\eqref{eq:avitzur-vega} therefore carries an explicit
assumption about the pressure on the land, $p = \sigma_f-\sigma_z$, which the exported contact data
can test.

The modification those authors propose is to drop the assumption that $\mu$ is a constant of the
lubricant and make it a function of the temperature measured by a thermocouple set in the die above
the bearing, which they take to be proportional to the temperature of the wire--die interface,
\begin{equation}
\mu = \mu_0\left(\frac{T}{T_0}\right)^{m_\mu},
\label{eq:mu-temperature}
\end{equation}
with $\mu_0$ the coefficient at a reference temperature $T_0$; inverting
Eq.~\eqref{eq:avitzur-vega} against their own force measurements they obtain $\mu_0 = 0.0125$,
$m_\mu = 2.68$ and $T_0 = 15.6\,^\circ$C for copper drawn in oil at 1--7~m/s, the temperatures
entering the ratio in degrees Celsius. In the present simulations $\mu$ is prescribed and constant,
so Eq.~\eqref{eq:mu-temperature} has no counterpart in the model; what it contributes to the
comparison is stated below.

All three relations are evaluated with the same representative flow stress derived from the
constitutive equation (Section~\ref{sec:prep-data}), 300--305~MPa for these runs, and with no
fitted parameter, so that the columns of Table~\ref{tab:drawstress} differ only in the form of
their work terms; the flow stress enters as a multiplier and shifts every prediction together.
% src: publication/reexport/verification_table.py (all columns); flow_stress.py (sigma_f);
% code/analytical.py (siebel, avitzur); code/vega_avitzur.py (Eq. avitzur-vega, Eq. mu-temperature)

\begin{table*}[htbp]
\centering\small
\caption{Drawing stress $\sigma_d$ (MPa) for the four runs: the finite-element value against the
Siebel relation, Eq.~\eqref{eq:siebel}, the Avitzur upper bound in its friction-factor form,
Eq.~\eqref{eq:avitzur} with $m=\sqrt3\,\mu$, and the same bound in the Coulomb form of
Eq.~\eqref{eq:avitzur-vega}. All columns use the flow stress derived from the constitutive
equation, 300--305~MPa, and no fitted parameter. The last column repeats
Eq.~\eqref{eq:avitzur-vega} with one change, the pressure on the land taken from the contact
reaction of the model, 122~MPa, instead of from the yield condition; it is a diagnostic of the
bearing assumption, not an independent prediction. The last row is the secant slope between the
$\mu = 0.025$ and $\mu = 0.1$ rows, so it carries the 1\,\% run-to-run change in the flow stress;
the finite-element value over the same two points is 323~MPa. Reduction 1.6\,\% in diameter
(3.2\,\% in area), land ratio unity, drawing speed 40~m/min.}
\label{tab:drawstress}
\begin{tabular}{@{}llrrrrr@{}}
\toprule
 & & & & & \multicolumn{2}{c}{Eq.~\eqref{eq:avitzur-vega}} \\
\cmidrule(l){6-7}
$\alpha$ & $\mu$ & FEM & Siebel & Avitzur & as published & land at 122~MPa \\
\midrule
$12^\circ$ & 0.025 & 54.1 & 53.2 & 75.1 & 71.5 & 66.1 \\
$12^\circ$ & 0.050 & 62.0 & 54.2 & 91.1 & 82.8 & 73.1 \\
$12^\circ$ & 0.100 & 78.3 & 57.4 & 125.4 & 104.4 & 88.8 \\
$8^\circ$  & 0.100 & 62.3 & 44.7 & 109.0 & 90.8 & 73.3 \\
\midrule
\multicolumn{3}{@{}l}{mean absolute deviation} & 17.3\,\% & 55.2\,\% & 36.2\,\% & 17.8\,\% \\
\multicolumn{3}{@{}l}{$\mathrm{d}\sigma_d/\mathrm{d}\mu$ at $\alpha=12^\circ$, MPa} & 57 & 670 & 439 & 303 \\
\bottomrule
\end{tabular}
\end{table*}
% src: publication/reexport/verification_table.py (columns, deviations, slope row, bearing
% pressures, sticking check); FEM row from drawing_force.py (time mean of |RF2| over die travel
% 20-100 mm, divided by the measured outlet section); sigma_f 302.4/301.1/305.3/299.9 MPa from
% flow_stress.py.

Both Avitzur columns lie above the computed drawing stress at every point, as the upper-bound
character of the underlying solution leads one to expect, although the theorem is not strictly
available here, since the model is elastic--plastic with hardening and thermal softening and
Eq.~\eqref{eq:avitzur-vega} is in addition a linearisation. The margins are not small, and they
order themselves by how each relation treats the land: 36.2\,\% on average in the Coulomb form and
55.2\,\% in the friction-factor form. Siebel is the closest, at 17.3\,\%, and it is the only one of
the three with no land term at all; it is also the only one that falls below the computed value,
which it does at every point, and its shortfall grows monotonically with friction, from 1.7\,\% at
$\mu = 0.025$ to 26.7\,\% at $\mu = 0.1$. Its low mean deviation is therefore a cancellation at the
lowest-friction run, not a uniform agreement. At a land ratio of unity the land is 89--93\,\% of
the contact area, so this is where the comparison is decided.
% src: verification_table.py (contact areas: land 1970 mm^2 against cone 156 mm^2 at 12 deg,
% 236 mm^2 at 8 deg)

The exported contact reaction locates the error. Through its yield closure $p=\sigma_f-\sigma_z$,
Eq.~\eqref{eq:avitzur-vega} implies a pressure of 240--251~MPa at the entrance to the land, falling
to a mean of 221--237~MPa along it, whereas the mean contact pressure in the model is 122~MPa: a
factor of 1.97 to 2.06 on the entrance value and 1.80 to 1.94 on the mean. The friction-factor form
is further out still, since $m=\sqrt3\,\mu$ is equivalent to taking the pressure on the land equal
to the flow stress, 300--305~MPa, a factor of 2.5. The sensitivity to friction confirms the
diagnosis independently of the magnitudes: from
$\mu = 0.025$ to $\mu = 0.1$ at $\alpha = 12^\circ$ the computed drawing stress rises by 323~MPa
per unit $\mu$, against 439 for Eq.~\eqref{eq:avitzur-vega}, 670 for Eq.~\eqref{eq:avitzur} and 57
for Siebel, whose insensitivity to friction is the other face of its missing land term.
Substituting the computed contact pressure into the bearing term and changing nothing else brings
the mean deviation to 17.8\,\% and the slope to 303~MPa per unit $\mu$, 6.1\,\% below the
322.7~MPa of the model.
The exports say directly that the yield closure does not hold there: on the land the measured
pressure and the axial stress sum to 176--200~MPa against a flow stress of 300--305~MPa, so the
material on the bearing is not at yield, and the two relations that assume it is overstate the only
term that carries most of their answer. That, and not the calibration of the flow stress, is what
both bearing models get wrong. The Coulomb branch itself remains admissible throughout: $\mu p$
reaches 25~MPa at most even at the pressure the relation assumes, and 12~MPa at the computed
pressure, against a shear yield stress of 173--176~MPa, so nowhere does the contact approach
sticking.
% src: verification_table.py (bearing pressure assumed at the land entrance and averaged along it,
% against the measured contact pressure; shear stress against the shear yield stress)

A second departure, of a different kind, survives that correction at the lowest friction. Setting
$\mu = 0$ in Eq.~\eqref{eq:avitzur-vega} leaves 58.8~MPa at $\alpha = 12^\circ$, already 8.8\,\%
above the 54.1~MPa the model returns with friction, so no treatment of the contact can close that
row at the flow stress used here. The offset lies in the redundant-shear term, which carries 29 to
65\,\% of the predicted drawing stress against 8 to 13\,\% for the ideal work. The inhomogeneity
parameter is 17.1 at
$\alpha = 8^\circ$ and 25.6 at $\alpha = 12^\circ$, against 1.81 in the case for which
Eq.~\eqref{eq:avitzur-vega} was calibrated, where the ideal work carries 74\,\% of the frictionless
drawing stress. The two failures are therefore distinct: the bearing term is wrong about the
contact, and the shear term is being used far outside the geometry in which the assumed velocity
field is a good approximation.
% src: verification_table.py (frictionless Eq. (1); term shares; Delta for the present runs and for
% the case of Vega et al.: R_i = 2.62e-4 m, R_f = 2.25e-4 m, alpha = 7.9 deg)

Equation~\eqref{eq:mu-temperature} cannot be tested on this campaign and must not be carried over
with its published constants. The Coulomb coefficient of the simulations is an input and does not
respond to the 15--21~$^\circ$C mean rise the model produces; and reading $\mu_0$ as the prescribed
coefficient at the initial 20~$^\circ$C would multiply it by 2.3 to 3.1 at the measured pass-mean
temperature of 27.5--30.6~$^\circ$C, or by 4.6 to 6.1 at the reference temperature of the original
calibration, in either
case moving every prediction further from the computed value rather than closer. What the relation
contributes here is the statement of what a friction model for this process would have to contain,
and the identification of what a campaign in which $\mu$ is prescribed rather than observed cannot
supply.
% src: verification_table.py (pass-mean temperature from the measured rise, 27.5-30.6 C;
% multipliers (T/20)^2.68 = 2.35-3.14 and (T/15.6)^2.68 = 4.56-6.11)

The comparison can be closed completely, with no quantity taken from the finite-element model. The
homogeneous strain is $2\ln(R_0/R_f)$ and the total strain follows from an excess-work
factor $\Phi = 0.8 + \Delta/4.4$; the strain rate is obtained from the time to traverse the cone at
constant volume flow, $\dot\varepsilon = 6\,v\,R_f^2\tan\alpha\,\ln(R_0/R_f)/(R_0^3-R_f^3)$; and the
flow stress is the work-weighted mean of the constitutive equation over that path:
\begin{align}
\varepsilon_{\mathrm{hom}} &= 2\ln\frac{R_0}{R_f}, \qquad
\Phi = 0.8 + \frac{\Delta}{4.4}, \qquad
\varepsilon_{\mathrm{eff}} = \Phi\,\varepsilon_{\mathrm{hom}}, \label{eq:closed-strain}\\
\dot\varepsilon &= \frac{6\,v\,R_f^2\tan\alpha\,\ln(R_0/R_f)}{R_0^3-R_f^3}, \qquad
\langle\sigma_f\rangle = A + \frac{B\,\varepsilon_{\mathrm{eff}}^{\,n}}{n+1}, \qquad
\Delta T = \frac{\beta\,\sigma_d}{\rho c}, \label{eq:closed-rest}
\end{align}
where $\Delta$ is the inhomogeneity parameter of the deformation zone and
$\langle\sigma_f\rangle$ the plastic-work-weighted mean of the strain term of the constitutive
equation, taken before its rate and temperature multipliers \citep{backofen1972deformation}. The
excess-work correlation is quoted in the literature in more than one form; replacing
$\Phi = 0.8+\Delta/4.4$ by the alternative $\Phi = 0.88+0.12\,\Delta$ changes $\Phi$ by a factor of
1.6 to 1.7 but the flow stress, and the estimate below, by only 1.3 to 1.6\,\%, because at
$\varepsilon_{\mathrm{hom}} = 0.032$ the hardening term of the constitutive equation is flat. The
choice therefore does not carry the comparison. The temperature
does not need a separate source, because it closes on the deformation work itself: that work per
unit volume is the part of $\sigma_d$ raised in the cone, a fraction $\beta$ of which becomes
heat, so the mean rise is $\Delta T = \beta\,\sigma_d/(\rho c)$; since $\sigma_d$ depends on the
flow stress and the flow stress on the temperature, this is a fixed point, reached in five
iterations. The friction on the land is left out of this balance deliberately: its heat is released
in a thin surface layer and is shared with the die, so it does not raise the mean temperature of
the section, whereas the friction on the cone, 2 to 14\,\% of the estimate, is retained because it
is dissipated in the deforming material itself. Carrying the land term into the balance as well
would raise the estimated rise by 25 to 120\,\% and break the agreement with the measured values of
Table~\ref{tab:deltaT}. For the four runs the fixed point returns a flow stress of 330--338~MPa.
This is 10\,\% above the 300--305~MPa used in Table~\ref{tab:drawstress}, and the difference is not
a disagreement: the fixed point evaluates the constitutive equation once, at
$\varepsilon_{\mathrm{eff}} = 0.215$ and 0.152, $\dot\varepsilon = 15.6$ and 10.3~s$^{-1}$ and the
pass-mean temperature, whereas the representative value of Section~\ref{sec:prep-data} averages the
same equation along the deformation path of every material point, most of which lies at lower
strain. The closed estimate is quoted with its own value throughout, and no column of
Table~\ref{tab:drawstress} uses it.
% src: publication/reexport/analytic_sigma_f.py (--Q 0.0161 --v 40; alpha 12/8, mu 0.025/0.05/0.1):
% eps_hom = 0.032, Phi = 6.62 / 4.70, eps_eff = 0.215 (alpha 12) / 0.152 (alpha 8),
% edot = 15.6 / 10.3 s^-1, sigma_f = 338.2/338.1/337.9/330.1 MPa. The land term is excluded from
% this closure (L/Rf = 0 in the solver), which is now stated in the text rather than left implicit.

\begin{table*}[htbp]
\centering\small
\caption{Intermediate quantities of the closed self-consistent estimate,
Eqs.~(\ref{eq:closed-strain})--(\ref{eq:closed-rest}),
for the four runs: inhomogeneity parameter $\Delta$, Backofen excess-work factor $\Phi$, effective
strain $\varepsilon_{\mathrm{eff}}$, mean strain rate $\dot\varepsilon$, flow stress $\sigma_f$ from
the constitutive equation at the fixed-point temperature, and the deformation part of the drawing
stress, $\sigma_d^{\mathrm{def}}$, being the ideal work, the redundant shear and the friction on the
cone, with the friction on the land excluded for the reason given in the text. Nothing is taken
from the finite-element model. The nominal geometry is used throughout,
$\varepsilon_{\mathrm{hom}}=0.032$ at $R_f=17.71$~mm, so $\Delta$ here is the value for that
geometry; evaluated instead at the outlet radius measured in each run it is 25.2--25.7 and 17.0.}
\label{tab:closed}
\begin{tabular}{@{}llrrrrrr@{}}
\toprule
$\alpha$ & $\mu$ & $\Delta$ & $\Phi$ & $\varepsilon_{\mathrm{eff}}$ & $\dot\varepsilon$, s$^{-1}$ & $\sigma_f$, MPa & $\sigma_d^{\mathrm{def}}$, MPa \\
\midrule
$12^\circ$ & 0.025 & 25.6 & 6.62 & 0.215 & 15.6 & 338.2 & 67.1 \\
$12^\circ$ & 0.050 & 25.6 & 6.62 & 0.215 & 15.6 & 338.1 & 68.4 \\
$12^\circ$ & 0.100 & 25.6 & 6.62 & 0.215 & 15.6 & 337.9 & 70.9 \\
$8^\circ$  & 0.100 & 17.1 & 4.70 & 0.152 & 10.3 & 330.1 & 53.9 \\
\bottomrule
\end{tabular}
\end{table*}
% src: publication/reexport/analytic_sigma_f.py --Q 0.0161 --v 40 (alpha 12/8, mu 0.025/0.05/0.1).
% The column is the deformation part only; the full drawing stress with the land is the
% "as published" column of Table tab:drawstress.

The temperature closure can be checked directly, since the dedicated export also carries the
temperature history. The self-consistent estimate reproduces the cross-section-averaged rise
measured in the model to within 12\,\% on average, and to within 19\,\% on every run
(Table~\ref{tab:deltaT}); the heating can therefore be taken
from the deformation work itself, with this comparison as its justification, rather than from an
assumed thermal boundary condition. The closure is not exact in its friction dependence: over the
friction range it moves the estimated rise by 1.0~$^\circ$C against 6.3~$^\circ$C measured, so some
of the land-friction heat does reach the section, but including that term in full would overshoot
the measured rise by more than it now falls short. Temperatures are stated in degrees Celsius
throughout, the
scale in which the constitutive model is defined, with a transition temperature of
20~$^\circ$C equal to the initial temperature of every run; a rise in degrees Celsius and a rise in
kelvin are the same number, but the absolute values are not, and
Eq.~\eqref{eq:mu-temperature} is calibrated on the Celsius scale.

\begin{table}[htbp]
\centering\small
\caption{Mean temperature rise during the pass and the resulting pass-mean temperature: the
self-consistent analytical estimate, in which the deformation work and the temperature close on
each other and nothing is read from the model, against the cross-section average of the exported
temperature history. Every run starts at 20~$^\circ$C, and the pass-mean temperature is formed from
the measured rise. The peak rise, reached at the surface, is 45.7--74.8~$^\circ$C, three to four
times the section mean. Reduction 1.6\,\% in diameter (3.2\,\% in area), land ratio unity.}
\label{tab:deltaT}
\begin{tabular}{@{}llrrr@{}}
\toprule
$\alpha$ & $\mu$ & $\Delta T$ analytical, $^\circ$C & $\Delta T$ measured, $^\circ$C & $T$ pass mean, $^\circ$C \\
\midrule
$12^\circ$ & 0.025 & 16.3 & 15.0 & 27.5 \\
$12^\circ$ & 0.050 & 16.6 & 16.9 & 28.4 \\
$12^\circ$ & 0.100 & 17.3 & 21.3 & 30.6 \\
$8^\circ$  & 0.100 & 13.1 & 15.8 & 27.9 \\
\bottomrule
\end{tabular}
\end{table}
% src: publication/reexport/analytic_sigma_f.py --Q 0.0161 (analytical);
% publication/reexport/drawing_force.py (measured: nodal peaks over the history, averaged over the
% radius with the area weight; the pass mean is 20 + dT_measured/2).

Three properties of the exported histories support the extraction independently of the analytical
comparison. The friction coefficient is recovered from the data rather than read from the file
name: over a die travel of 115 to 140~mm the cone has left the wire and only the cylindrical land
remains in contact, the contact normal is then strictly radial, and the ratio of the axial to the
radial reaction returns 0.0250, 0.0500, 0.1000 and 0.0999 for the four runs. The mean contact
pressure is 122~MPa and varies by 1~MPa over the four runs, with neither the semi-angle nor the
friction, as it must when a cylindrical land is 89--93\,\% of the contact area. And the flow stress recovered from
the heat alone, over the inner half of the radius where friction contributes none, is 292--329~MPa,
in agreement with the 300--305~MPa derived from the constitutive equation; beyond that radius the
same estimate runs up past 500~MPa as the friction heat of the land reaches it, which is itself a
check that the core is the right place to take it.
% src: publication/reexport/drawing_force.py (mu from data, contact pressure over the steady
% window). The recovered coefficient also settles the identity of the fourth run, whose file name
% carries a different value: the export is one run at mu = 0.1 (ERRATA E-23).

The region of the design covered by these runs limits what the comparison establishes. All four
lie at a reduction of 1.6\,\% with a land as long as the outlet radius, a combination in which the
inhomogeneity parameter is 25.6 at $\alpha = 12^\circ$ and 17.1 at $\alpha = 8^\circ$, against the
range of 1 to 3 in which these solutions are normally applied and against which they have been
validated; they remain valid outside it but become progressively loose, because the assumed
velocity field departs further from the actual flow. The ideal work accounts for 8 to 13\,\% of the
predicted drawing stress and the total exceeds it by a factor of 7.7 to 12.4.
The deviations of Table~\ref{tab:drawstress} therefore characterise the closed-form solutions in
this corner of the design, not the accuracy of the finite-element model, and they do not transfer
to the industrial range of reductions, where the ideal work dominates and the land is a small
fraction of the contact. What does transfer is the bearing diagnosis: wherever the land is long
relative to the reduction, a closed-form solution that puts it at the yield pressure will
over-predict the drawing force by about as much as that pressure is overstated.

In addition, five internal consistency tests were run on the 1756 retained cases. The exported cross-section is self-equilibrated to 3.2\,\%: the median $|\int\sigma_{zz}\,\mathrm{d}A|/A$ is 5.3~MPa against an rms $\sigma_{zz}$ of 167~MPa. % src: VALIDATION.md §7 (5.3 / 167 → 3.2 %)
Local equilibrium closes to 2.1\,\% of the sum of the moduli of its terms. % src: VALIDATION.md §7 (2.1 %)
At the free surface $|\sigma_{rr}|$ is 4.0~MPa against 54.0~MPa in the bulk, i.e.\ 7.4\,\% (a nodal value; the 3.52~MPa quoted below is the median of the outermost bin of the 20-bin profile), % src: VALIDATION.md §7 (4.0 / 54.0 → 7.4 %); RESULTS.md header (3.52 MPa on the binned profile); TODO: confirm the node population and averaging behind the 4.0 MPa of VALIDATION.md §7 (the file does not state them)
on the axis $|\sigma_{rr}-\sigma_{\theta\theta}| = 1.07$~MPa (0.6\,\%), % src: VALIDATION.md §7 (1.07 MPa → 0.6 %)
and the residual field stays below the flow stress everywhere: the maximum von Mises stress over the retained cases is 232.6~MPa, against the 300--305~MPa that the constitutive equation returns for the deformation histories of this pass (Section~\ref{sec:prep-data}).
% REV: the comparison is now made against the flow stress derived from the constitutive equation,
% not against a postulated yield stress of 320 MPa. % src: VALIDATION.md §7 (max Mises 232.6 at σ_y = 320)
Kinematically, the exit radius follows $R_0(1-Q)$ to 0.11\,\% (ratio 1.0011) and $|\mathrm{tr}\,\mathbf{LE}| = 5.1\times10^{-4}$, i.e.\ plastic incompressibility holds to 0.05\,\%. % src: VALIDATION.md §7 (1.0011 → 0.11 %; 5.1e-4 → 0.05 %)
These tests check the mesh, the contact and the solver of this model, not its correspondence to real wire. % src: VALIDATION.md §5 item 6
